% Part: proof-theory
% Chapter: sequent-calculus
% Section: introduction

\documentclass[../../../include/open-logic-section]{subfiles}

\begin{document}

\olfileid[pt]{pt}{seq}{int}

\olsection{Introdução}

Os cálculos de sequentes são sistemas de !!{proof} introduzidos na década
de 1930 por Gerhard Gentzen. Eles não foram concebidos propriamente como um
método prático para !!{proving} coisas. Gentzen os usou, antes, para poder
demonstrar certos resultados sobre a possível estrutura dos !!{proof}s e
suas propriedades. Questões desse tipo são precisamente aquelas de que se
ocupa a teoria da prova.

Os cálculos de sequentes operam sobre \emph{sequentes}, como sugere o nome,
e não sobre !!{formula}s. Um sequente é um par de sequências ou
multiconjuntos de !!{formula}s. Os sistemas originais de Gentzen (\Log{LK}
e \Log{LJ}) usavam sequências; hoje, os teóricos da prova preferem usar
multiconjuntos. Um multiconjunto é semelhante a um conjunto, exceto por
poder conter um !!{element} mais de uma vez. Contudo, assim como nos
conjuntos, a ordem dos !!{element}s não importa. Portanto, as expressões
$\{!A, !A, !B\}$ e $\{!A, !B, !A\}$ descrevem o mesmo multiconjunto, mas
esse multiconjunto é diferente de $\{!A, !B\}$ e de $\{!A, !A, !B, !B\}$,
pois os últimos contêm números de $!A$ e $!B$ diferentes dos do primeiro.

Tradicionalmente, multiconjuntos de !!{formula}s na teoria da prova são
representados por letras gregas maiúsculas. Assim, quando escrevemos
$\Gamma$, $\Delta$, $\Pi$, $\Lambda$ ou $\Theta$, pretendemos que sejam
multiconjuntos de !!{formula}s da lógica com a qual estivermos trabalhando.
Também por tradição, em vez de escrever $\tuple{\Gamma, \Delta}$, os
teóricos da prova separam os dois componentes por um símbolo de sequente;
nós usamos~$\Sequent$. 

\begin{defn}[Sequent]
Um \emph{sequente} é uma expressão da forma
\[
\Gamma \Sequent \Delta
\]
onde $\Gamma$ e $\Delta$ são multiconjuntos finitos (possivelmente vazios)
de !!{formula}s. $\Gamma$ é chamado de \emph{antecedente}, enquanto
$\Delta$ é chamado de \emph{sucedente}.
\end{defn}

Seja $!G$ a conjunção das !!{formula}s em~$\Gamma$ (tomando a conjunção
vazia como~$\ltrue$), e seja $!D$ a disjunção das !!{formula}s em~$\Delta$
(tomando a disjunção vazia como~$\lfalse$). Então, o sequente
$\Gamma \Sequent \Delta$ é verdadeiro (em alguma estrutura~$\Struct{M}$)
se, e somente se, $!G \lif !D$ é verdadeiro. Na lógica clássica, isso
equivale a dizer que alguma das !!{formula}s em~$\Gamma$ é falsa ou alguma
das !!{formula}s em~$\Delta$ é verdadeira.

Se $\Gamma$ é um multiconjunto de !!{formula}s, escrevemos $\Gamma, !A$
(ou $!A, \Gamma$) para o resultado de acrescentar $!A$ a~$\Gamma$. Se
$\Delta$ também é um multiconjunto de !!{formula}s, então $\Gamma, \Delta$
é a união multiconjuntista das duas sequências---se $\Gamma$ contém $!A$
com \emph{multiplicidade}~n (O.s.\@, contém $n$ cópias de~$!A$) e $\Delta$
com multiplicidade~$m$, então $\Gamma, \Delta$ contém $!A$ com
multiplicidade~$n+m$. Se o multiconjunto em um dos lados de um sequente é
vazio, em geral simplesmente deixamos um espaço em branco. Por exemplo,
$\Sequent !A$ é o sequente com antecedente vazio e cujo sucedente contém
$!A$ com multiplicidade~$1$.

!!^a{proof} no cálculo de sequentes é uma árvore de sequentes, na qual os
sequentes no topo (ou iniciais) são axiomas do sistema considerado e, se um
sequente aparece abaixo de um ou dois outros sequentes, deve decorrer deles
corretamente por uma regra de inferência do sistema. Consideraremos primeiro
dois sistemas de sequentes diferentes para a lógica clássica, $\Log{G1c}$
e~$\Log{G3c}$. Seus axiomas e regras são apresentados em
\cref{pt:seq:int:tab:G1c,pt:seq:int:tab:G3c}. \oliflabeldef{fol:seq::chap}{ (O
sistema original de Gentzen \Log{LK} é descrito em \olref[fol][seq][]{chap}.)}{}

\subfile{rules-G1c}
\subfile{rules-G3c}

Os sistemas diferem de maneiras sutis, e essas diferenças fazem com que
tenham propriedades distintas do ponto de vista da teoria da prova.
\Log{G1c} tem axiomas $!A \Sequent !A$ nos quais nenhuma outra das
!!{formula}s é permitida. \Log{G3c} tem axiomas da forma
$!C, \Gamma \Sequent \Delta, !C$, com !!{formula}s ``laterais'' arbitrárias
em $\Gamma$ e~$\Delta$, mas a !!{formula}~$!C$ que aparece dos dois lados
deve ser \emph{atômica}. \Log{G1c} tem duas versões de cada uma das regras
$\LeftR{\land}$ e $\RightR{\lor}$, enquanto \Log{G3c} tem apenas uma de
cada. Além disso, \Log{G1c} tem ``regras estruturais'' adicionais chamadas
de ``contração'' (\LeftR{\Contraction}, \RightR{\Contraction}) e
``enfraquecimento'' (\LeftR{\Weakening}, \RightR{\Weakening}).

Os sistemas \Log{G1c} e \Log{G3c} são corretos e completos para a lógica
clássica. Em outras palavras, todo sequente !!{provable} nesses sistemas é
verdadeiro em toda !!{structure}, e todo sequente verdadeiro em toda
!!{structure} tem !!a{proof}. Assim, a questão de \emph{se} um sequente tem
!!a{proof} também pode ser respondida por outros métodos da lógica (p.ex., os da teoria de modelos). E, como ambos os sistemas !!{prove}
exatamente os sequentes verdadeiros em toda !!{structure}, em particular
eles demonstram os mesmos sequentes.

Contudo, a teoria da prova não se ocupa apenas de \emph{se} algo é
demonstrável, mas de \emph{como} o é. Os teóricos da prova consideram
questões como ``É sempre possível evitar certa regra?'', ``Essa regra pode
ser acrescentada ao sistema sem tornar demonstráveis novos sequentes?'' ou
``É possível transformar !!{proof}s de um sistema em !!{proof}s de outro?''
Se a resposta a uma pergunta desse tipo for ``sim'', esse fato costuma ser
demonstrado por indução sobre a complexidade dos !!{proof}s ou,
equivalentemente, por indução sobre a estrutura dos !!{proof}s. Isso fornece
não apenas uma demonstração do fato (p.ex., se \Log{G1c} demonstra um
sequente, então \Log{G3c} demonstra o mesmo sequente), mas também um
procedimento (p.ex., para transformar !!{proof}s em~\Log{G1c} em
!!{proof}s em~\Log{G3c}).

Um resultado central da teoria da prova é o \emph{teorema da eliminação do
corte}. Ele mostra que podemos acrescentar a regra do corte
\[
\Axiom$\Gamma \fCenter \Delta, !A$
\Axiom$!A, \Pi \fCenter \Lambda$
\RightLabel{\Cut}
\BinaryInf$\Gamma, \Pi \fCenter \Delta, \Lambda$
\DisplayProof
\]
aos sistemas de sequentes sem alterar quais sequentes são !!{provable}.
Além disso, a demonstração fornece um método que transforma qualquer
!!{proof} que use as regras de \Log{G1c} (ou \Log{G3c}) juntamente com
\Cut{} em outro que usa somente as regras de \Log{G1c} (ou \Log{G3c}).
Em outras palavras, esse procedimento ``elimina'' a regra \Cut{} dos
!!{proof}s que a usam; daí o nome.

\end{document}
